\documentclass[10pt,a4paper]{amsart}
\usepackage[margin=19mm]{geometry}
\usepackage{amsmath,amssymb,mathtools}
\usepackage[scr=rsfs,cal=euler]{mathalfa}
\usepackage{xfrac}
\usepackage[unicode,hidelinks,bookmarks=false]{hyperref}
\newcommand{\ie}{i.e.\ }
\newcommand{\cf}{cf.\ }
\newcommand{\resp}{respectively\ }
\newcommand{\Subsection}{\subsection}
\providecommand{\nobreakdash}{\nobreak\hbox{-}}
\setlength{\emergencystretch}{2em}
\multlinegap0pt
\usepackage{bm}
\usepackage{bbm}
\usepackage{dsfont}
\usepackage{xspace}
\usepackage{cancel}
\usepackage{mathtools}
\usepackage{cleveref}
\usepackage[shortlabels]{enumitem}
\usepackage{amsthm}
\makeatletter
\@ifundefined{Theorem}{\newtheorem{Theorem}{Theorem}}{}
\@ifundefined{Corollary}{\newtheorem{Corollary}[Theorem]{Corollary}}{}
\@ifundefined{Definition}{\newtheorem{Definition}[Theorem]{Definition}}{}
\@ifundefined{Lemma}{\newtheorem{Lemma}[Theorem]{Lemma}}{}
\@ifundefined{Notation}{\newtheorem{Notation}[Theorem]{Notation}}{}
\@ifundefined{Proposition}{\newtheorem{Proposition}[Theorem]{Proposition}}{}
\@ifundefined{Remark}{\newtheorem{Remark}[Theorem]{Remark}}{}
\makeatother
\def\KK{{\mathbb K}}
\def\ini{\operatorname{in}} 
\newcommand\bdX{{\bm X}}
\newcommand\bda{{\bm a}}
\newcommand\bdc{{\bm c}}
\newcommand\bfS{\mathbb{S}}
\newcommand\bfX{\mathbf{X}}
\newcommand\calC{\mathcal{C}}
\newcommand\calF{\mathcal{F}}
\newcommand\calL{\mathcal{L}}
\newcommand{\Lr}{\ensuremath{\calL\times[r]}}
\title[Gorenstein Special Fiber Rings of Ladder Determinantal Modul]{Gorenstein Special Fiber Rings of Ladder Determinantal Modules}
\author{Louiza Fouli, Kuei-Nuan Lin, Haydee Lindo, Maral Mostafazadehfard}
\date{Source: arXiv:2601.13483v1 (CC BY 4.0)}
\begin{document}
\maketitle

\noindent\emph{Source and licence.} Gorenstein Special Fiber Rings of Ladder Determinantal Modules. Version arXiv:2601.13483v1 (CC BY 4.0), distributed under the Creative Commons Attribution 4.0 International licence, as recorded in the arXiv licence metadata for this exact version and shown on the abstract page https://arxiv.org/abs/2601.13483v1. Full rights notice: licenses/arxiv-2601.13483-CC-BY-4.0.txt.\par\smallskip

\noindent\textit{Editorial scope.} Counted unit: the Proposition whose source label is \texttt{components} in the article (source0.tex lines 882--885, source label \texttt{components}), delivered together with the article's own complete original proof of it (source0.tex lines 887--904, one \texttt{proof} environment of 18 lines). No mathematics is written, repaired or re-derived: statement and proof are the article's own text line for line. Required context retained uncounted: def:Ladder (lines 654--670); shrinkmatrix (lines 726--733); cor: join irr (lines 800--802); cor: north east (lines 813--815); lem: coverinP (lines 838--842); not: disconnected partition (lines 863--875). Every definition, display and label of those lines is retained unchanged. Omitted from this package and not redistributed: 1--653, 671--725, 734--799, 803--812, 816--837, 843--862, 876--881, 886--886, 905--1320. Nothing inside a retained excerpt is omitted or altered: every retained body is delivered exactly as the article's lines are written, so no omission receipt of any kind accompanies this package. Citations: the retained excerpts carry no citation command at all, so no bibliography entry of the article is transcribed and no reference is redistributed. Reference adaptation: none was needed and none was made; every reference inside the delivered unit resolves inside it, so no unresolved reference or citation is produced. Printed slips kept literal: no printed word, symbol, hypothesis or number of the counted statement or of its proof was altered. Layout and class: the article's own class is \texttt{amsart} with packages including geometry, amsfonts, amsmath, amssymb, stmaryrd, amsthm, amsxtra, bbm, braket, comment, extarrows, graphicx, latexsym, mathrsfs; the wrapper is the lane's pinned \texttt{amsart} editorial wrapper at 10pt on A4, which restates the theorem-like environments the retained text needs and adds the article's own macro definitions that the retained text needs (\texttt{\textbackslash KK}, \texttt{\textbackslash Lr}, \texttt{\textbackslash bdX}, \texttt{\textbackslash bda}, \texttt{\textbackslash bdc}, \texttt{\textbackslash bfS}, \texttt{\textbackslash bfX}, \texttt{\textbackslash calC}, \texttt{\textbackslash calF}, \texttt{\textbackslash calL}, \texttt{\textbackslash ini}), with extra wrapper packages (bm, bbm, dsfont, xspace, cancel, mathtools, cleveref, enumitem). The delivered numbering is the unit's own. Assets: no retained span contains a figure environment, a graphics-inclusion command, a PDF-page inclusion, a table or a TikZ picture; no image file is copied into this package, the package declares no asset dependency and no image file is redistributed with it. Licence: version arXiv:2601.13483v1 of this article is distributed under the Creative Commons Attribution 4.0 International licence, as recorded in the arXiv licence metadata for this exact version and shown on the abstract page https://arxiv.org/abs/2601.13483v1. A full rights notice is delivered at licenses/arxiv-2601.13483-CC-BY-4.0.txt. Characters: every retained excerpt is pure ASCII, so the delivered engine, XeLaTeX with its default Latin Modern text font, reports no missing character.
\begin{Definition}
    \label{def:Ladder} 
    Let $\KK$ be a field and let $m,n,r$ be positive integers with $n<m$.  Write $[m]=S_1 \cup S_2 \cup \cdots \cup S_n$, where for each $i\in [n]$, $S_i\coloneqq [u_i, v_i]$ for some $u_i, v_i$ positive integers with $u_i < v_i$.  Without loss of generality, we may assume that $1=u_1<  u_2 <  \cdots <  u_n < m$ and $1< v_1<  v_2 < \cdots < v_n=m$.
     For each $i\in [n]$ and each $j\in S_i$, let $x_{i,j}$ be an indeterminate over the field $\KK$. Let $\bfS \coloneq \{S_1, \ldots, S_n\}$.


    \begin{enumerate}[a]
   \item    Let  $\bfX_{\bfS}$ be the $n\times m$ matrix whose $(i,j)$-th entry is $x_{i,j}$ when $j\in S_i$ and $0$ otherwise.  We say that $\bfX_{\bfS}$ is the \emph{ladder matrix} associated to $\bfS$. 
    \item Let $R\coloneq\KK[\bfX_{\bfS}]=\KK[x_{i,j}: i\in[n], j\in S_i]$ be a polynomial ring over $\KK$ and $L\coloneq I_n(\bfX_{\bfS})$ be the $R$-ideal generated by the $n\times n$ (maximal) minors of $\bfX_{\bfS}$. 
     \item Let  $\calL\coloneqq \{\bdc=(c_1, \ldots, c_n)\in S_1\times \cdots \times S_n : c_i<c_{i+1} \mbox{ for all } i\in [n-1]\}$. Notice that by the ladder shape of $\bdX_{\bfS}$ we have  $\det[\bdc]\neq 0$ for any $\bdc\in \calL$, where  $\det[\bdc]$ is the determinant of the $n\times n$ submatrix of $\bfX_{\bfS}$ whose columns are $c_1, \ldots, c_n$.  Therefore, the set $\{ \det[\bdc] : \bdc\in \calL\}$ is a minimal generating set of the ideal $L$.
     
     \item The module $M=\bigoplus_{i=1}^rL$ is called a \emph{ladder determinantal module} associated to the matrix $\bfX_{\bfS}$. Let $\calL\times [r]  \coloneqq \{(\bdc, j) : \bdc\in \calL, j\in [r]\}$ be the cartesian product of $\calL$ and $[r]$ and notice that the set $\{(\det[\bdc], j): (\bdc, j)\in \Lr\}$ is a minimal generating set for $M$.
     \item Let $\tau$ be the lexicographic order on the ring $R$ induced by the lexicographic order on the entries in $\bfX_{\bfS}$, that is $x_{i,j}>_{\tau}x_{l,k}$ if and only if $i<l$ or $i=l$ and $j<k$.
     \item Let $N=\bigoplus_{i=1}^r \ini_{\tau}(L)$ be the direct sum of the initial ideal of $L$ with respect to $\tau$.
     \item For all $i\in[n]$ and all $j\in [n-1]$, let $\Delta_i\coloneq v_i-u_i$, $\epsilon_j\coloneq u_{j+1}-u_j$, and  $\theta_j\coloneq v_{j+1}-v_j$. We adopt the convention that $\epsilon_n = \theta_0=\infty$.
\end{enumerate}
\end{Definition}

\begin{Remark}\label{shrinkmatrix}
For all $i \in [2,n]$, let $S_i' = S_i$ if $u_{i-1} < u_i$, and let $S_i' = [u_i + 1, v_i]$ whenever $u_{i-1} = u_i$.  
Let $\bfS' = \{S_1, S_2', \ldots, S_n'\}$, and let ${\bfX}_{\bfS'}$ be the ladder matrix corresponding to $\bfS'$.  
Notice that $\ini_{\tau}(I_n(\bfX_{\bfS})) = \ini_{\tau}(I_n(\bfX_{\bfS'}))$.  
Furthermore, since $\tau$ is a diagonal monomial order, we may assume that $u_i \le v_{i-1} + 1$ for all $i \in [2,n]$.  
Hence, 
we may assume without loss of generality that $u_i < u_{i+1}$ and $v_i < v_{i+1}$ as in  \Cref{def:Ladder}. If $u_i=v_i$ for some $i\in [n]$, then $L=x_{i,u_i}I_{n-1}(\bfX'_{\bfS})$, where $\bfX'_{\bfS}$ is the ladder matrix obtained from $\bfX_{\bfS}$ after removing row $i$. Note that $\calF(M) \cong \calF(\bigoplus_{i=1}^r I_{n-1}(\bfX'_{\bfS})$ and thus we may assume without loss of generality that $u_i<v_i$ for all $i\in [n]$ as in \Cref{def:Ladder}. Moreover, under our assumptions we cover the case of a generic matrix by considering $\bfX_{\bfS}$ with $\epsilon_i=1$ and $\theta_i=1$ for all $i\in [n-1]$. Finally, if $n=m$, then $\calF(M)$ is a polynomial ring and there is nothing to prove in this case. Therefore we may assume that $n<m$. 
\end{Remark}

\begin{Corollary}\label{cor: join irr}
    The join-irreducible elements of $\calL\times [r]$ are precisely the elements \[ P_{\Lr} := \{(\bda_{i,j},1): i\in[n], j\in[u_i+1, v_i]\}\cup \{\bda_{n+1,k}: k\in [2,r]\}.\]
\end{Corollary}

\begin{Lemma} \label{cor: north east}
       Let $\bda_{i,j}, \bda_{k,l}\in P_{\calL}$. Then $\bda_{k,l}\le \bda_{i,j}$ if and only if $k\ge i$ and $l-j\le k-i$.
\end{Lemma} 

\begin{Lemma} \label{lem: coverinP}
If $\bda_{i,j}\in P_{\calL}$, then the only elements $\bda_{i,j}$ covers are the elements in the set $$ \{\bda_{i,j-1}, \bda_{i+1,j+1}\},$$ which may be the empty set.
Moreover, 
$\bda_{i,j-1} \in P_{\calL} $ if and only if $ j \ge u_i+2$ and $\bda_{i+1, j+1} \in P_{\calL} $ if and only if $j \ge u_{i+1}$ and $i\neq n$.  
\end{Lemma}

\begin{Notation}\label{not: disconnected partition}
    Consider the set $$\calC=\{i\in [n-1]: v_i<u_{i+1}\} \cup \{n\}=\{q_1, \ldots, q_{t-1}, q_t=n: q_{i}<q_{i+1} \text{ for all } i\in[n-1]\}.$$ 
Notice that $i\in \calC$ if and only if $u_{i+1}=v_i+1$ for all $i\in [n-1]$ by \Cref{shrinkmatrix}.
    We partition $[n]$ as follows \[[n]=[1, q_{1}]\cup [p_{2}, q_2]\cup \cdots \cup[p_{t}, q_t],\] where $p_{i+1}=q_{i}+1$ for all $i\in[t-1]$. 
    For simplicity let $p_1=1$.
Notice that with these assumptions $\bfX_{\bfS}$ is a block matrix with $t$ blocks $\bfX_i$, where $\bfX_i$ is a ladder matrix in the rows $p_i$ to $q_i$:
    $$\bfX_{\bfS}=\begin{bmatrix}
        \bfX_1 &&&\\
        &\bfX_2&&\\
        &&\ddots&\\
        &&&\bfX_t\\
    \end{bmatrix}.$$
\end{Notation}

\begin{Proposition} \label{components}
 The lattice $P_{\Lr}$ is connected if and only if $r=1$ and $|\calC|=1$.
   In particular, $P_{\Lr}$ has $|\calC|$ connected components when $r=1$ and $|\calC| +1$ connected components when $r> 1$. 
\end{Proposition}

\begin{proof}
Recall that  $P_{\Lr} = \{(\bda_{i,j},1): i\in[n],  j\in[u_i+1, v_i]\}\cup \{\bda_{n+1,k}: k\in [2,r]\}$, \Cref{cor: join irr}.  Let $\calC=\{i\in [n-1]: v_i<u_{i+1}\} \cup \{n\}=\{q_1, \ldots, q_{t-1}, q_t=n: q_{i}<q_{i+1} \text{ for all } i\in[n-1]\}$ be as in \Cref{not: disconnected partition}. For every $s\in [t]$ let $A_s\coloneq\{(\bda_{i,j},1): i\in [p_s,q_s], j\in[u_i+1,v_i]\}$. 
Let $B=\{\bda_{n+1,k}:k\in[2,r]\}$, which exists if and only if $r\ge 2$. We claim that $A_1, \ldots, A_t, B$ are the connected components of $P_{\Lr}$ when $r\ge 2$. 
In fact, $A_s$ is an isomorphic copy of the poset corresponding to the $s$-th block of $X_{\mathbb{S}}$.

Note that $B\cup \bigcup_{i=1}^t A_i=P_{\Lr}$. 
Observe that $(\bda_{i,j},1)$ is incomparable to $\bda_{n+1,k}$ for any $i\in [n]$, $j\in [u_i+1,v_i]$ and $k\in [2,r]$. Moreover, $B$ is  a connected component of $P_{\Lr}$ since $\bda_{n+1,k}<\bda_{n+1,k+1}$ for all $k\in [2,r-1]$.

 Let $s\in [t]$ and let $i\in [p_s,q_s]$. By \Cref{cor: north east} the set $\{(\bda_{i,j},1): j\in [u_i+1,v_i]\}$ is connected. If $i<q_s$, then  
$v_i\ge u_{i+1}$. Hence $(\bda_{i,v_i},1)>(\bda_{i+1,v_i+1},1)$ by \Cref{lem: coverinP}. Therefore, $A_s$ is connected.


 

We claim that for any $s, s'\in [t]$ with $s\neq s'$ the elements of $A_s$ and $A_{s'}$ are incomparable. Indeed, let $(\bda_{i,j},1)\in A_s$ and $(\bda_{g,h},1)\in A_{s'}$ and without loss of generality assume $s'>s$. Then $g>i$ and $\bda_{g,h} \not> \bda_{i,j}$. Since 
$$j\le v_i\le v_{q_s}<u_{q_s+1}=u_{p_{s+1}} \le u_g<h,$$ 
it follows that $h-j>u_g-v_i\ge g-p_{s+1}+1+q_s-i=g-i$, as $p_{s+1}=q_s+1$, $u_{k}-u_l\ge k-l$ and $v_k-v_l\ge k-l$ for any $k>l$. Thus,  $\bda_{g,h}\not<\bda_{i,j}$ and hence $ (\bda_{i,j},1), (\bda_{g,h},1)$ are incomparable by \Cref{cor: north east} as claimed. Therefore, $A_1, \ldots, A_t, B$ are the connected components of $P_{\Lr}$. In the case $r=1$, one can show as above that the sets $A'_s=\{\bda_{i,j}: i\in [p_s, q_s] ,j\in [u_i+1,v_i] \}$ are the connected components of $P_{\calL}$ for all $s\in [t]$.
\end{proof}

\end{document}
